\chapter{Limiti tecnici}
\label{chap:conclusions}

\paragraph{Riassunto del capitolo.} Prima di concludere, riconosciamo
sei limiti reali del lavoro svolto: l'incapacità di estrapolare oltre
il bordo dell'\emph{Aging} osservato, l'assenza di modellazione della
variabilità inter-cella, la mancanza di feedback per il concept drift,
il \emph{tuning paradox} dovuto a fold di CV informativamente
eterogenei, l'assenza di intervalli di confidenza calcolati per
bootstrap sulle metriche e la copertura zero del frontend. Sono i punti
da aggredire per primi se il lavoro venisse esteso a un secondo
progetto.

\section{Estrapolazione fuori dal dominio osservato}

Il dataset misura cinque livelli di Aging in una finestra
\emph{specifica}. Predire $Aging = 5$ (oltre il bordo) non è supportato
dal nostro setup attuale: KNN ricadrebbe sui vicini più prossimi (i
quattro Aging già osservati) e il valore predetto sarebbe
sistematicamente \emph{sotto-stimato}. Una mitigazione sarebbe usare
modelli con prior fisico (es.~Arrhenius parametrizzato sulla
temperatura, polinomio di basso grado in Aging) per estendere la base.

\section{Cross-cell variability}

Il dataset deriva da \emph{una sola} cella pouch \acrshort{licoo2} da
\SI{10}{\ampere\hour}. La variabilità inter-cella (tolleranze di
produzione, batch, condizioni di formazione) non è modellata. Per usare
la pipeline su un altro form-factor o un'altra chimica (LFP, NMC)
servirebbe un nuovo training; il modulo di drift detection suonerebbe
l'allarme appena le distribuzioni di $\mathrm{Re}(Z), \mathrm{Im}(Z)$
uscissero dall'inviluppo del dataset attuale.

\section{Concept drift}

Il modulo di monitoring rileva \emph{covariate shift}, non
\emph{concept drift}. In mancanza di un canale di feedback con
\emph{label veri} per le predizioni in produzione, non è possibile
osservare il deterioramento del modello in modo diretto. In un sistema
di laboratorio una soluzione realistica sarebbe \emph{re-misurare
periodicamente un sottoinsieme di celle} e confrontare le predizioni
con le misure dirette.

\section{Tuning paradox}

In tutti e tre i task il GridSearchCV non produce un miglioramento sul
fold di test rispetto ai default. La formulazione corretta è
statistica: la valutazione finale è single-shot (un solo hold-out
fold), mentre la selezione del grid è fatta su una media \acrshort{cv}
di 4--39 fold \acrshort{logo} informativamente molto eterogenei. Senza
un intervallo di confidenza (bootstrap, cfr.~\cref{sec:bootstrap}) non
è lecito concludere che ``il tuning peggiora'': è più verosimile che il
valore tunato cada \emph{dentro} il CI del default.

I sintomi residui sono comunque coerenti con tre ingredienti noti:
\begin{enumerate}[leftmargin=*]
  \item dataset relativamente piccolo (\num{9805} righe $=$ $\sim$200
        curve);
  \item griglia di iperparametri larga rispetto alla varianza
        intra-fold;
  \item splitter LOGO con fold di bordo informativamente diversi.
\end{enumerate}

\section{Bootstrap dei CI sulle metriche}
\label{sec:bootstrap}

Nessuna delle tabelle riportate nei \cref{chap:task1,chap:task2,chap:task3}
include intervalli di confidenza sulle metriche. È un limite reale: una
differenza di MSE non è interpretabile senza una stima di varianza. La
soluzione standard è il bootstrap di Efron~\cite{efron1979bootstrap}:
ricampionare il fold di test con sostituzione $B \approx 1000$ volte e
riportare il quantile 2.5\,\% e 97.5\,\% di ciascuna metrica. È
un'aggiunta di poche righe in \codefile{src/evaluation/metrics.py}.

\section{Frontend testing}

La copertura di test sul frontend è zero. Per un progetto di semestre
è accettabile (UI minimale, nessun stato globale complesso); per andare
in produzione bisognerebbe aggiungere test di componente
(React Testing Library, Jest) e un E2E (Playwright o Cypress) per
verificare che l'integrazione frontend--backend non si rompa.
